\documentclass[../../../include/open-logic-section]{subfiles}

\begin{document}

\olfileid{his}{set}{infinitesimals}
\olsection{بی‌نهایت‌کوچک‌ها و مشتق‌گیری}

نیوتن و لایبنیتس در پایان سدهٔ هفدهم، مستقل از یکدیگر، حساب دیفرانسیل و
انتگرال را کشف کردند. یکی از کاربردهای بسیار مهم آن \emph{مشتق‌گیری}
بود. مشتق‌گیری، به‌تقریب، می‌کوشد مفهومی از «نرخ تغییر» یا شیب تابع در یک
نقطه به دست دهد.

روش روشنی برای نمایش این ایده چنین است. تابع
$f(x) = \nicefrac{x^2}{4} + \nicefrac{1}{2}$ را در نظر بگیرید که در شکل
زیر با رنگ سیاه رسم شده است:
\begin{center}
	\begin{tikzpicture}[scale=1]
	\draw[->, gray] (-1,0) -- (4.25,0) node[right] {$x$};
	\draw[->, gray] (0,-0.25) -- (0,5.25) node[above] {$f(x)$};

	\foreach \x/\xtext in {1/1, 2/2, 3/3, 4/4}
	\draw[shift={(\x,0)}, gray] (0pt,2pt) -- (0pt,-2pt) node[below] {$\xtext$};

	\foreach \y/\ytext in {1/1, 2/2, 3/3, 4/4, 5/5}
	\draw[shift={(0,\y)}, gray] (2pt,0pt) -- (-2pt,0pt) node[left] {$\ytext$};

	\draw[oldiagcolorC] (0.5, 9/16) -- (3.5, 57/16) -- (3.5, 9/16)--cycle;
	\draw[oldiagcolorD] (.5, 9/16) -- (2.5, 33/16) -- (2.5, 9/16) -- cycle;
	\draw[oldiagcolorE] (.5, 9/16) -- (1.5, 17/16) -- (1.5, 9/16) -- cycle;
	\draw[thick] (-1,.75) parabola bend (0,0.5) (4,4.5);
	\end{tikzpicture}
\end{center}
فرض کنید می‌خواهیم شیب تابع را در $c = \nicefrac{1}{2}$ بیابیم. کار را
با رسم مثلثی آغاز می‌کنیم که وتر آن شیب تابع در آن نقطه را تقریب می‌زند؛
مثلاً مثلث قرمز بالا. اگر $\beta$ طول قاعدهٔ مثلث باشد، ارتفاع آن
$f(\nicefrac{1}{2}+\beta) - f(\nicefrac{1}{2})$ است؛ پس شیب وتر برابر
است با:
\[
\frac{f(\nicefrac{1}{2}+\beta) - f(\nicefrac{1}{2})}{\beta}.
\]
بنابراین شیب مثلث !!{colorC} ما با طول قاعدهٔ~$3$ دقیقاً~$1$ است. وتر
مثلثی کوچک‌تر، یعنی مثلث !!{colorD} با طول قاعدهٔ~$2$، تقریب بهتری به دست
می‌دهد و شیبش $\nicefrac{3}{4}$ است. مثلثی باز هم کوچک‌تر، یعنی مثلث
!!{colorE} با طول قاعدهٔ~$1$، تقریب بهتری می‌دهد و شیبش
$\nicefrac{1}{2}$ است.

مثلث‌های هرچه کوچک‌تر تقریب‌های هرچه بهتری به دست می‌دهند. پس ممکن است
بگوییم وتر مثلثی با طول قاعدهٔ \emph{بی‌نهایت‌کوچک} خودِ شیب تابع را در
$c = \nicefrac{1}{2}$ به دست می‌دهد. به این ترتیب، فرمول مشتق (مرتبهٔ
اول) تابع $f$ در نقطهٔ $c$ چنین خواهد بود:
\[
{f'}(c) = \frac{f(c+\beta) - f(c)}{\beta}
\text{ که در آن $\beta$ بی‌نهایت‌کوچک است.}
\]
به‌تقریب، نیوتن و لایبنیتس همین را گفتند.

اما چون آنان چنین گفتند، باید بپرسیم: یک \emph{بی‌نهایت‌کوچک} چیست؟
دوراهی دشواری پیش می‌آید. اگر $\beta = 0$ باشد، $f'$ تعریف‌نشده است،
زیرا تقسیم بر $0$ در آن رخ می‌دهد. اما اگر $\beta > 0$ باشد، تنها یک
\emph{تقریب} از شیب به دست می‌آوریم، نه خود شیب را.

این نگرانی ناهم‌زمان با تاریخ نیست. برکلی در نقد پیروان نیوتن می‌نویسد:
\begin{quote}
می‌پذیرم که نشانه‌ها می‌توانند بر چیزی یا بر هیچ‌چیز دلالت کنند؛ و در
نتیجه در نمادگذاری آغازین $c + \beta$، ممکن بود $\beta$ بر افزایشی دلالت
کند یا بر هیچ‌چیز. اما هرکدام را که معنای آن قرار دهید، باید سازگار با
همان معنا استدلال کنید و بر معنایی دوگانه تکیه نکنید؛ زیرا چنین کاری
مغالطه‌ای آشکار است. (\citealt[\S{}XIII]{Berkeley1734}، متغیرها برای
هماهنگی با متن پیشین تغییر یافته‌اند)
\end{quote}
برای دفاع از حساب بی‌نهایت‌کوچک‌ها در برابر برکلی، ممکن است پاسخ دهیم که
سخن‌گفتن از «بی‌نهایت‌کوچک‌ها» صرفاً مجازی است. می‌توان گفت تا وقتی مثلثی
\emph{واقعاً کوچک} بگیریم، تقریبی \emph{به‌اندازهٔ کافی خوب} از خط مماس
خواهیم داشت. برکلی برای این سخن نیز پاسخی داشت: شاید این مقدار برای مهندسی
کافی باشد، اما \emph{جایگاه} ریاضیات را تضعیف می‌کند، زیرا
\begin{quote}
به ما گفته‌اند که \emph{in rebus mathematicis errores qu\`{a}m minimi
non sunt contemnendi}. [در ریاضیات، کوچک‌ترین خطاها را نباید نادیده
گرفت.] \citep[\S{}IX]{Berkeley1734}
\end{quote}
عبارت ایتالیک تقریباً نقل‌قولی لفظ‌به‌لفظ از \emph{تربیع منحنی‌ها}
(۱۷۰۴)، اثر خود نیوتن، است.

اعتراض‌های فلسفی برکلی بسیار ژرف و نافذند. بااین‌حال، حساب دیفرانسیل و
انتگرال طرحی بسیار موفق بود و ریاضی‌دانان بی‌آنکه دچار خطا شوند به
استفاده از آن ادامه دادند.

\end{document}
